\section{冻结 OPD 信号与 KL 方向}
\label{opd:mirror-origin}

固定 prefix $s$，本章写 $q=q_k(\cdot\mid s)$、$p=p_{k,n}(\cdot\mid s)$、$t=t(\cdot\mid s)$，三者在有限词表上严格为正。第一章命题~\ref{prop:foundation-kl-objective} 已从 teacher KL 的镜像线性化推导冻结信号 $r=\log(t/q)$。本章将同一结论代入具体信息访问与缓存契约；$r^{\mathrm{cur}}=\log(t/p)=r-U$ 是另一项随 student 改变的 current reward。

本章带分布方向的 KL 下标是全书旧中心记号的显式别名：
\begin{align}
K_{qT,k}(s)&=\mathcal K_k^T(s)=\KL(q_k(\cdot\mid s)\Vert t(\cdot\mid s)),\\
K_{qp,k,n}(s)&=\mathcal K_{k,n}^S(s)=\KL(q_k(\cdot\mid s)\Vert p_{k,n}(\cdot\mid s)),\\
K_{pT,k,n}(s)&=\mathcal K_{k,n}^T(s)=\KL(p_{k,n}(\cdot\mid s)\Vert t(\cdot\mid s)).
\end{align}
前者在 outer 冻结，后两者在 inner 重算。省略 $k,n,s$ 后写成 $K_{qT},K_{qp},K_{pT}$；$K_S$ 与 $K_T$ 只在明确说明采样测度的局部推导中使用。这里三个数来自同一分布三元组，并未引入额外 KL 目标。

\begin{corollary}[OPD 的镜像目标实例]
\label{opd:objective-family}
应用命题~\ref{prop:foundation-kl-objective}，本章的 MDPO 目标满足
\begin{equation}
F_\eta(p)=\eta^{-1}\KL(p\Vert q)-\E_p[r]
=\KL(p\Vert t)+(\eta^{-1}-1)\KL(p\Vert q).
\end{equation}
因此 $\eta=1$ 时 MDPO 与 teacher KL 在任意候选 student 上处处相同。
\end{corollary}

$\eta<1$ 对应 teacher KL 与正的 anchor 正则；$\eta=1$ 消除了总体目标的 anchor 依赖；$\eta>1$ 外推最优 log-ratio，其完整目标的概率坐标 Hessian 仍为 $\operatorname{diag}(1/(\eta p_a))$，故严格凸；$\eta\to\infty$ 则变为冻结 reward 的线性优化。后一极限与 current-reward OPD 的 teacher KL 目标不同。这些等式固定了 $s$ 和 $d_k$。若优化整条 response 分布的 KL，或让 prefix 占用测度随 student 改变，求导还要处理前缀概率与未来代价；不能把本章的条件 next-token 梯度直接解释成那个动态轨迹目标的完整梯度。

\begin{proposition}[$\eta=1$ 的 MDPO 与 Token-IS OPD]
\label{opd:eta-one-mdpo}
记 $\psi_a=\nabla_\theta\log p_a$、$U_a=\log(p_a/q_a)$、$r_a^{\mathrm{cur}}=\log(t_a/p_a)=r_a-U_a$。在同一固定 prefix 测度下，MDPO 与 current-reward Token-IS OPD 具有相同总体梯度；二者的有限样本系数可以相差 baseline。
\end{proposition}
\begin{proof}
当行为分布为 $q$ 时，MDPO 系数是 $(p_a/q_a)(U_a-r_a)$，OPD 系数是 $(p_a/q_a)(U_a-r_a-K_{pT})$。两者之差的期望为
\begin{align}
\E_q\left[\frac pqK_{pT}\psi\right]
&=K_{pT}\sum_a p_a\nabla_\theta\log p_a\\
&=K_{pT}\nabla_\theta\sum_a p_a=0.
\end{align}
任意覆盖 $p$ 的行为分布 $\beta$ 用 $p/\beta$ 得到同样结论。baseline 在取样之前给定或与主 action 条件独立时，这个零积分论证成立；若用同一 action 估计 baseline，还须核对自耦合项。有限冻结数据上的更新、不同 prefix 刷新和 optimizer 状态不由总体等价自动确定。
\end{proof}

\section{OPD 特化：逐 token reward 与两种信息访问模式}
\label{sec:ch2-opd-specialization}

在 OPD 中，第 $k$ 个 outer step 先固定 anchor $q_k$，然后从它产生 rollout。对 prefix $s_t=(x,y_{<t})$ 和实际 sampled token $a_t=y_t$，最直接的局部 teacher reward 是
\begin{equation}
\label{eq:ch2-opd-raw-reward}
r_k(s_t,a_t)=\log t(a_t\mid s_t)-\log q_k(a_t\mid s_t).
\end{equation}
它满足序列分解
\begin{equation}
\sum_{t=1}^H r_k(s_t,a_t)=\log t(y\mid x)-\log q_k(y\mid x).
\end{equation}
因此在共同 mirror 子问题中应取 $f=r_k$；不能把整条序列的 log-ratio 当作每一个 prefix 的同一个 reward，否则会重复计算后缀差异。

可以对 reward 做两种不同的“归一化”，必须区分它们的理论作用。对任意只依赖 prefix 的 baseline $b_k(s)$，令
\begin{equation}
\label{eq:ch2-opd-centered-reward}
f_k^{\mathrm{ctr}}(s,a)=r_k(s,a)-b_k(s).
\end{equation}
它只把 partition function 乘上 $e^{-\eta b_k(s)}$，所以 Gibbs policy 不变；取 $b_k(s)=\mathbb E_{q_k}[r_k(s,\cdot)\mid s]$ 就是条件中心化。若再除以尺度 $\sigma_k(s)>0$，则
\begin{equation}
f_k^{\mathrm{std}}=(r_k-b_k)/\sigma_k
\end{equation}
会把有效温度从 $\eta$ 改成 $\eta/\sigma_k$，不再是同一个 mirror step，除非同时把正则参数改为 $\eta\sigma_k$。因此 OPD 的默认实现只做 additive centering；标准化只能作为明确的 temperature schedule。

\subsection{完整 vocabulary logits：精确 token mirror target}
\label{sec:ch2-opd-full-vocab}

完整词表 logits 使冻结 reward $r_k$、teacher center $K_{qT,k}$、student center $K_{qp,k,n}$ 和 partition function 都可精确求和。OPD 的配分函数为
\begin{equation}
\log Z_k(s)=\operatorname{logsumexp}_{v\in\mathcal V}
\left((1-\eta)\log q_k(v\mid s)+\eta\log t(v\mid s)\right).
\label{eq:ch2-opd-full-logz}
\end{equation}
相应 Gibbs target 是
\begin{equation}
p_k^{\mathrm G}(v\mid s)
=\exp\!\left((1-\eta)\log q_k(v\mid s)+\eta\log t(v\mid s)-\log Z_k(s)\right).
\end{equation}
PMD-Mean 的 target 是定理~\ref{thm:ch2-mean-lambert} 的 $p_k^{\mathrm M}$。完整词表路径消除了条件 token 采样误差；prefix coverage 与数值优化误差仍保留。

\subsubsection*{四种损失与实际梯度}
第一章的条件均值恒等式给出 $\mathbb E_{q_k}[r_k\mid s]=-K_{qT,k}$、$\mathbb E_{q_k}[U_{k,n}\mid s]=-K_{qp,k,n}$。因此冻结 Mean target 为 $G_k=\eta(r_k+K_{qT,k})$，三个回归 residual 是
\begin{align}
R_{k,n}^{\mathrm{Part}}(s,v)&=U_{k,n}(s,v)-\eta r_k(s,v)+\log Z_k(s),
\label{eq:ch2-opd-full-part-simplified}\\
R_{k,n}^{\mathrm{Mean}}(s,v)&=U_{k,n}(s,v)-\eta(r_k(s,v)+K_{qT,k}(s)),
\label{eq:ch2-opd-full-mean-simplified}\\
R_{k,n}^{\mathrm{BPO}}(s,v)&=U_{k,n}(s,v)+K_{qp,k,n}(s)-\eta(r_k(s,v)+K_{qT,k}(s)).
\label{eq:ch2-opd-full-bpo-residual}
\end{align}
它们对应
\begin{align}
\mathcal L_{\mathrm{Part}}(s)&=\tfrac12\mathbb E_{q_k}[(R_{k,n}^{\mathrm{Part}})^2\mid s],\\
\mathcal L_{\mathrm{Mean}}(s)&=\tfrac12\mathbb E_{q_k}[(R_{k,n}^{\mathrm{Mean}})^2\mid s],\\
\mathcal L_{\mathrm{BPO}}(s)&=\tfrac1{2\eta}\mathbb E_{q_k}[(R_{k,n}^{\mathrm{BPO}})^2\mid s],
\label{eq:ch2-opd-full-bpo-simplified}\\
\mathcal L_{\mathrm{MDPO}}(s)&=\mathbb E_{p_{k,n}}[U_{k,n}/\eta-r_k\mid s].
\label{eq:ch2-opd-full-mdpo-simplified}
\end{align}
三个平方损失的解析梯度为
\begin{align}
\nabla_\theta\mathcal L_{\mathrm{Part}}(s)&=\mathbb E_{q_k}[R_{k,n}^{\mathrm{Part}}\psi_{k,n}\mid s],\\
\nabla_\theta\mathcal L_{\mathrm{Mean}}(s)&=\mathbb E_{q_k}[R_{k,n}^{\mathrm{Mean}}\psi_{k,n}\mid s],\\
\nabla_\theta\mathcal L_{\mathrm{BPO}}(s)&=\frac1\eta\mathbb E_{q_k}[R_{k,n}^{\mathrm{BPO}}(\psi_{k,n}-\bar\psi_{k,n})\mid s]\\
&=\mathbb E_{q_k}\left[\left(\frac{U_{k,n}+K_{qp,k,n}}\eta-r_k-K_{qT,k}\right)\psi_{k,n}\mid s\right].
\label{eq:ch2-opd-full-bpo-gradient-simplified}
\end{align}
这里直接比较 $\nabla\mathcal L$，BPO 的 $\eta^{-1}$ 已乘入权重。其实际下降方向为
\begin{equation}
-\nabla_\theta\mathcal L_{\mathrm{BPO}}(s)
=\mathbb E_{q_k}\left[\left(r_k+K_{qT,k}-\frac{U_{k,n}+K_{qp,k,n}}\eta\right)\psi_{k,n}\mid s\right].
\label{eq:ch2-opd-bpo-descent-expanded}
\end{equation}
centered score 的第二项消失，是因为 $\mathbb E_{q_k}[R_{k,n}^{\mathrm{BPO}}\mid s]=0$；不能把 $\bar\psi_{k,n}$ 自身视为零。Mean residual 的 anchor 均值为 $-K_{qp,k,n}$，所以 Mean 不具有同样的消去。

MDPO 同时微分 student 测度和 log-ratio，得到
\begin{align}
\nabla_\theta\mathcal L_{\mathrm{MDPO}}(s)
&=\mathbb E_{p_{k,n}}\left[\left(\frac{U_{k,n}+1}\eta-r_k\right)\psi_{k,n}\mid s\right]\\
&=\mathbb E_{p_{k,n}}[(U_{k,n}/\eta-r_k)\psi_{k,n}\mid s].
\end{align}
第二行使用 student 测度的零均值 score 恒等式。MDPO 使用 student 积分测度，三个回归损失使用 anchor 积分测度。

\subsubsection*{停止梯度的具体位置}
$q_k$、$t$、$\log Z_k$、$K_{qT,k}$ 和 $G_k$ 在固定 outer step 冻结。student KL 每个 inner step 重算，原始 BPO 平方损失中保留其导数：
\begin{align}
K_{qp,k,n}(s)&=\mathbb E_{q_k}[\sg(\log q_k)-\log p_{k,n}\mid s],\\
\nabla_\theta K_{qp,k,n}(s)&=-\bar\psi_{k,n}(s).
\label{eq:ch2-opd-grad-ktheta}
\end{align}
若实现已经解析好的 policy gradient，所有 residual 和 MDPO 权重都整体停止梯度：例如 surrogate
\begin{equation}
\mathbb E_{q_k}[\sg(R_{k,n}^{\mathrm{Mean}})\log p_{k,n}]
\end{equation}
的梯度正好等于 Mean 梯度。此处权重的导数已由解析推导处理，再微分一次会改变目标。若直接实现原始平方损失，则 student residual 及 BPO center 保留梯度。Mean 原始损失不含 student KL center。
